% Replace the former exact-compilation subsection with this fragment.
% Retain the preceding shared-point arithmetic source and its hypotheses.
% No claim here removes or proves its simultaneous-prime input.
\subsection{The quadratic compilation in 962 variables}
\label{subsec:v6-962-compilation}

We compile the shared-point source above using different point charts,
pole coordinates, and arithmetic circuits. Every change has explicit
rational witness maps. The arithmetic source, including its simultaneous
prime input, is unchanged. Throughout this subsection
\[
 K=\Q(\sqrt2),\qquad D=(3-\sqrt2)^2=11-6\sqrt2,
 \qquad \mathsf S=3^{2^{27}}.
\]
An equation in $K$ means its two rational component equations. The fixed
integer $\mathsf S$ is a coefficient, not a coordinate. An assignment of
a quadratic expression to one rational coordinate adds one quadratic
equation. Affine expressions require no coordinate.

\begin{theorem}[The 962-variable quartic]
\label{thm:v6-962-polynomial}
The complete table specified below defines
\[
 P_{962}\in\Z[X_0,\ldots,X_{961}]
\]
with degree four, exactly $55409$ nonzero coordinate monomials, integer
coefficient content one, and coefficient height
\[
 H(P_{962})=64\mathsf S^2=64\cdot3^{2^{28}}<2^{429496736}.
\]
It is the sum of squares of $919$ quadratic residuals and has the same
rational input projection as the shared-point source of
Theorem~\ref{thm:v5-source-projection}. Consequently, under that theorem's
arithmetic hypotheses, for every $x\in\Q$,
\[
 x\in\Z\quad\Longleftrightarrow\quad
 \exists X_1,\ldots,X_{961}\in\Q,\quad
 P_{962}(x,X_1,\ldots,X_{961})=0.
\]
The $962$ variables include the input and $961$ existential witnesses.
\end{theorem}

\paragraph{A quadratic chart for the nine points.}
For each point introduce $t,h\in K$ and set
\begin{equation}\label{eq:v6-point-chart}
 u=t^2,\qquad X=h^2,\qquad
 u(X-u/4)=D,\qquad Y=h(u/2-X).
\end{equation}
The three assignments use six rational coordinates. The displayed
constraint contributes two quadratic residuals. The four rational
components of $t,h$ replace the four halfpoint components.

\begin{lemma}[Exact point-chart maps]\label{lem:v6-point-chart}
The chart \eqref{eq:v6-point-chart} describes precisely the finite doubles
used in the shared-point source, with the same rational witnesses up to
the following invertible changes. It preserves all required nonzero
point-coordinate domains. The rational component $X_0$ of $X$ is positive.
\end{lemma}
\begin{proof}
For an old halfpoint $(a,b)$ with $b^2=a^3-Da$ and its tangent equation,
both $a,b$ are nonzero. Indeed $b=0$ would give $a=0$ or $a^2=D$,
whereas the tangent numerator $3a^2-D$ would be $-D$ or $2D$. Then
$a=0$ is excluded by the curve equation. Define
\[
 t=b/a,\qquad h=(a^2+D)/(2b)=a/t-t/2.
\]
The curve equation gives $t^2=a-D/a$. Substitution gives
$t^2(h^2-t^2/4)=D$. The doubling formulas give exactly $X=h^2$
and $Y=h(t^2/2-h^2)$.

Conversely, a chart solution has $t\ne0$ because $D\ne0$. Set
\[
 a=th+t^2/2,\qquad b=at,\qquad m=h+t.
\]
If $a=0$, then $h=-t/2$, contradicting the chart equation. Thus $b\ne0$.
Direct expansion gives
\[
 a^2-t^2a=D,\quad b^2=a^3-Da,\quad
 2bm=3a^2-D,\quad m^2-2a=h^2,
\]
and $m(a-h^2)-b=h(t^2/2-h^2)$. The substitutions are inverse, so they
recover exactly the old finite double. Lemma~\ref{lem:v5-point-domains}
therefore supplies its nonzero abscissa and ordinate, and the same remains
true for the unchanged guarded differences and subsequent doublings.

Both real embeddings of $D$ are positive. Since $t\ne0$, the chart
equation implies $h^2=t^2/4+D/t^2>0$ under both embeddings. Hence
$X_0$, the average of these two positive values, is positive.
\end{proof}

The five affine subtraction formulas keep their original inverses of
the differences of abscissas. Neither a subtraction guard nor a point
occurrence is deleted. Relative to the direct halfpoint doubling circuit,
the chart saves two rational arithmetic coordinates for each of the nine
points.

\paragraph{The logarithm polynomial.}
Retain the source's coefficients
\[
 a_0=1,\quad a_1=-2/5,\quad
 a_j=a_{j-1}(4j-5)/(4j+1).
\]
For $P(w)=\sum_{j=0}^{12}a_jw^j$ write
\begin{equation}\label{eq:v6-log-square-completion}
 P(w)=a_{12}Q(w)^2+R(w),\qquad
 \deg Q=6,\quad Q\text{ monic},\quad\deg R\le5.
\end{equation}
This identity is constructed without any witness division. Starting
with the leading coefficient of $Q$ equal to one, match degrees eleven
down to six successively. At each step the new coefficient is determined
by division by $2a_{12}\ne0$. The remainder then has degree at most five.
After computing $w^2,w^3$, every term through $w^6$ is quadratic in
$1,w,w^2,w^3$. Assigning $Q$ and then $a_{12}Q^2+R$ therefore uses
four field assignments, rather than five. The eight order-$52$
logarithms save sixteen rational coordinates. The order-$16$ logarithms
keep their exact degree-three polynomials. No truncation order changes.

Usually the final logarithm is $\ell=zP(Df^2)$, with
$f=1/x$ and $z=-x/y$. For $U_0$ and $V_0$ only, omit the separate $z$
coordinates and impose the direct quotient equations
\begin{equation}\label{eq:v6-log-quotient}
 y\ell=-xP(Df^2).
\end{equation}
The polynomial output is an affine expression in assigned coordinates,
so these are quadratic component equations. The ordinate is nonzero by
Lemma~\ref{lem:v5-point-domains}; hence \eqref{eq:v6-log-quotient}
is exactly the old logarithm. The source's logarithm-nonvanishing result
and every logarithm-ratio domain remain unchanged. In the scalar block
replace $z(H)\rho=z(S)$ by
\begin{equation}\label{eq:v6-scalar-ratio}
 x(H)\rho=-y(H)z(S).
\end{equation}
Both $x(H),y(H)$ are nonzero, so this has the same unique value
$\rho=z(S)/z(H)$.

\paragraph{Products, tests, and detector evaluation.}
For each of the three orders, compute the elementary symmetric
coefficients of $\prod_j(1+qD_j^2)$ by a balanced product tree.
At an internal node with coefficient vectors $a_i,b_j$, assign each
component of $\sum_{i+j=k}a_ib_j$ as one quadratic expression.
This is the ordinary coefficient-convolution identity. It preserves
every coefficient and every test, including its assigned order.

The degree-ten detector polynomial is evaluated in two blocks:
\begin{equation}\label{eq:v6-detector-block-evaluation}
 \sum_{j=0}^{10}c_jq^j
 =\sum_{j=0}^{4}c_jq^j
   +q^5\left(\sum_{j=0}^{5}c_{j+5}q^j\right).
\end{equation}
Compute $q^2,q^3,q^4,q^5$ using four scalar assignments, assign the
two components of the parenthesized block, and assign the two final
components directly. The low block is not assigned separately.
Every final summand is quadratic in the available coordinates. This
uses eight rational coordinates. The degree-two and degree-four
polynomials use their direct power sums, with their two component
sums assigned once. These formulas are identities over $K$.

The Lucas expression is regrouped as
\[
 t^2-a(a-2)tb+((a-2)b)^2-1
 =t^2-(at)C+C^2-1,\qquad C=(a-2)b.
\]
Thus one field multiplication is removed without changing the tested
polynomial. The fixed quintic is evaluated by its exact coefficients
using $t^2$; it remains
$F_5(t)=t(t+1/2)(t+1)(t-1/2399)(t+32/407)$.

\paragraph{Weight quotients and their domains.}
Fix a detector and write $A$ for its scale, $\delta=A\lambda-1$,
and $e$ for its weight. The scales are
\[
 1,\quad x(Q)_0^2,\quad (1+x(Q)_1^2)^2
 \qquad(Q=A,S,T,Z).
\]
They are all positive by Lemma~\ref{lem:v6-point-chart}. Retain the
inverse $\delta i=1$, and replace the separate fifth-power product
and inverse of $F_5$ by
\begin{equation}\label{eq:v6-weight-domain}
 F_5(A\lambda)e=(A\delta)^4(A\delta).
\end{equation}
Here $(A\delta)^4$ denotes the output of its retained power-assignment
gates; the displayed product is quadratic in the available coordinates.
The right side is nonzero, so this equation forces both $F_5(A\lambda)$
and $e$ nonzero. In particular the removed inverse of $F_5$ is recovered
uniquely. Define the low-order weight by
\[
 \delta q_1=e i^2,
\]
where $i^2$ is likewise a retained square-assignment output. This gives
the unchanged $q_1=e/\delta^3$. Retain the repeated-squaring
values $e^6,e^{24},e^{96}$ and set
\[
 c=q_1e^6,\qquad q_{16}=ce^{24},\qquad q_{52}=ce^{96}.
\]
These are the old $q_1e^{30}$ and $q_1e^{102}$, with one shared product.

If the detector product is $W=w_0+w_1\sqrt2$, retain $w_0j=1$ and use
\begin{equation}\label{eq:v6-norm-final-weight}
 w_0n=(w_0^2-2w_1^2)/2,\qquad ev=in,\qquad B=nv.
\end{equation}
The trace guard is essential and is retained. It gives $W\ne0$ in $K$,
so $\operatorname{Norm}(W)\ne0$ and
$n=\operatorname{Norm}(W)/(2w_0)\ne0$. Equations
\eqref{eq:v6-weight-domain}--\eqref{eq:v6-norm-final-weight} therefore
give
\[
 B=\frac{i}{e}\left(\frac{\operatorname{Norm}(W)}{2w_0}\right)^2\ne0,
\]
the original detector input. The reverse map restores $R=i/e$, its
old inverse intermediates, and $n^2$ by their unique rational values.
No inverse is inferred from an unguarded zero-over-zero quotient.

\paragraph{The scalar norm and finite corrections.}
The scalar contact polynomial $r$ and its inverse $ri=1$ remain.
Writing $f=f(H)_0$, $u=f(H)_1$, and $g=f(S)_0$, retain the leaf
$\ell=ui$ and replace two norm-tree
nodes by
\[
 rn_0=r+f^2,\qquad rn_1=u\ell+g^2.
\]
with the existing leaf $ui$ retained. Since $r\ne0$, these equal
$n_0=1+r(fi)^2$ and $n_1=(ui)^2+r(gi)^2$. All other nodes and all
orders in the scalar norm tree are unchanged.

Replace the eight Boolean correction choices and their seven product
coordinates by $c_0,\ldots,c_7\in\Q$ satisfying
\begin{equation}\label{eq:v6-correction-chain}
 c_0^2=1,\qquad
 (c_j-c_{j-1})(c_j-p_jc_{j-1})=0\quad(1\le j\le7),
\end{equation}
where $(p_1,\ldots,p_7)=(2,3,5,7,11,13,17)$. Induction shows
$c_j\ne0$ and $c_j/c_{j-1}\in\{1,p_j\}$. Thus $c_7$ is exactly an
allowed signed squarefree correction. Conversely every original
choice determines such a chain. Recover the original Booleans by
\[
 \epsilon_-=(1-c_0)/2,\qquad
 \epsilon_j=\frac{c_j/c_{j-1}-1}{p_j-1}.
\]
All denominators are nonzero. The scalar preparation coordinate,
its defining quotient, and the explicit inverse of its pole input
are retained.

\paragraph{The ten pole blocks.}
Write $b$ for a pole input, $C=17b$, and $U$ for its prepared value.
For the nine coordinate poles $U=b$. In the scalar pole $U$ retains
the source's preparation formula and its denominator domain.
Every $C$ is nonzero: this was proved above for coordinate detectors,
and the scalar input has its retained inverse.

Replace the old witness $e_0$ by $z=Ce_0$. Introduce the five
intermediates
\[
 P=p^2,\quad F=U^2,\quad H=z(C-F),\quad V=s^2,\quad Q=q^2,
\]
and impose
\begin{equation}\label{eq:v6-scaled-poles}
 CP=(C-z)(C+3z),\qquad CQ=(V-H)(V+3H).
\end{equation}
Every equation is quadratic. Under $z=Ce_0$, the first residual is
$C$ times the old residual $p^2-C(1-e_0)(1+3e_0)$, and the second
is unchanged because $H=Ce_0(C-U^2)$. The inverse map $e_0=z/C$
is defined. This saves one arithmetic coordinate at each pole,
without changing the pole predicate or any preparation condition.

\paragraph{Counts and the final polynomial.}
The initial allocation is
\[
\begin{array}{l|r|r}
 &\text{coordinates}&\text{structural residuals}\\\hline
 \text{input}&1&0\\
 \text{nine point charts}&36&18\\
 \text{common slope}&1&0\\
 \text{nine coordinate poles}&36&18\\
 \text{scalar pole with preparation}&5&2\\
 \text{private anchor}&3&1\\
 \text{signed correction chain}&8&8\\\hline
 \text{total}&90&47
\end{array}
\]
The complete program has $904$ arithmetic coordinates and their
$904$ defining equations. Hence the raw counts are $994$ coordinates
and $951$ quadratic residuals.

The private anchor uses exactly the old affine change
\[
 u=\mathsf S+(U-3V+3W)/4,\quad
 v=\mathsf S/2+V/4,\quad w=3\mathsf S/2+W/4.
\]
Its determinant is $1/64$. If $T=f_0^2+f_0f_1-f_1^2$, sixteen times
the old anchor after substitution is
\[
 -48T^2+U^2-6UV+6UW+15V^2-18VW+7W^2+8\mathsf S U.
\]
No anchor condition or witness is removed.

Each certified compiler elimination solves a row $aY+q=0$ with
$a\in\Q^*$ and $q$ independent of $Y$, substitutes $Y=-q/a$
everywhere, and removes that row. It is accepted only if every
remaining residual is still quadratic. The reverse map reinserts
the unique displayed value of $Y$. There are exactly $32$ such
eliminations, no row additions, and no omission of an unused coordinate.
The input is retained. Consequently
\[
 90+904-32=962,\qquad47+904-32=919.
\]
Normalize the remaining rows to primitive integer polynomials
$q_1,\ldots,q_{919}$ and define
\begin{equation}\label{eq:v6-final-polynomial}
 P_{962}(X_0,\ldots,X_{961})=\sum_{\nu=1}^{919}q_\nu^2.
\end{equation}
The complete coefficient table is
\nolinkurl{quartics/quartic-962/polynomial.json}, with SHA-256
\begin{center}\small
\nolinkurl{6b798de7d00ac3c7d725ca0a69d39f1eac1ddb02fee2756cb6323770b952af32}.
\end{center}
An entry with coefficient $c$, exponent \texttt{sPower} $k$, and
coordinate list $[i_1,\ldots,i_d]$ means
$c\mathsf S^kX_{i_1}\cdots X_{i_d}$. Every entry is included and repeated
indices give powers. This defines an ordinary integer polynomial.

\begin{proof}[Proof of Theorem~\ref{thm:v6-962-polynomial}]
For a source solution apply the point-chart, correction-chain and
pole substitutions, and evaluate all new assignment and quotient
coordinates in dependency order. The identities and domain proofs
above show that every structural residual vanishes. Conversely apply
the inverse substitutions, restore every deleted quotient by its
proved nonzero denominator, and evaluate the source assignments.
Every source point, logarithm, test, detector input, scalar value and
anchor value is recovered. These maps fix $X_0$. The affine anchor
change, certified eliminations and nonzero row normalizations preserve
the same projection. Over $\Q$, \eqref{eq:v6-final-polynomial} vanishes
exactly when all $919$ residuals vanish. This proves the projection
statement, and Theorem~\ref{thm:v5-source-projection} gives its conditional
arithmetic interpretation.

Independent exact convolution gives $55409$ distinct monomials.
Of their coefficients, $55401$ are $\mathsf S$-independent integers,
seven are $c\mathsf S$ with $0<|c|\le768$, and one is
$64\mathsf S^2$. The largest absolute $\mathsf S$-independent coefficient
has bit length $1526$. Thus no coefficient vanishes on specialization
and $64\mathsf S^2$ is the exact height. It occurs at $X_{79}^2$,
where $X_{79}=U$. The coefficient of $X_0^4$ is one, proving degree
four and integer content one. Finally
$\log_2(64\mathsf S^2)<6+(8/5)2^{28}<429496736$ because $3^5<2^8$.
\end{proof}

The frozen package includes the builder, raw systems, all $904$ gate
definitions, the exact elimination and normalization certificate,
compact coordinate map, final residuals, complete coefficient table,
source-equivalence proofs and independent reviews. Its manifest has
SHA-256
\begin{center}\small
\nolinkurl{36fca72cb0eb126193ffa43d890ff09d5e29cbf1274561d2632ed1e227163787}.
\end{center}
The independent checker verifies every gate, all $894$ certificate
operations, the complete sum-of-squares expansion, support,
integrality, content and height. A fresh isolated reproduction gives
the same seven deterministic artifacts byte for byte. These finite
checks establish the stated algebraic facts; they do not prove any
inherited arithmetic input.
